\subsection{Operations of a groupoid}

Let G be a groupoid; denote by S the set of its vertices. Let X be a set and let \(p\colon \mathrm{X}\to \mathrm{S}\) be a map. A (right) operation \(\varphi\) of the groupoid G on the set X, relative to \(p\) , is a morphism of groupoids \(\varphi\) from G to the groupoid \(\mathcal{B}(\mathrm{X},p)\) of the permutations of the set X relative to \(p\) (II, p. 163) which induces the identity on the set of vertices. In other words, \(\varphi\) is the data, for every pair \((a,b)\) of points of S and for every arrow \(g\in \mathrm{G}\) joining \(a\) to \(b\), of a map \(\varphi (g)\colon \mathrm{X}_a\to \mathrm{X}_b\) , such that for every pair \((f,g)\) of composable arrows of G,

\[
\varphi (f g) = \varphi (g) \circ \varphi (f),
\]

and such that, for every \(a \in S\) , \(\varphi(e_a)\) is the identity of \(X_a\) .

Let \(\varphi\) and \(\varphi'\) be operations of the groupoid G on a set X, relative to a map \(p\colon X\to S\) . For \(\varphi=\varphi'\) , it suffices that there exist a subquiver H of G, generating G, such that \(\varphi(f)=\varphi'(f)\) for all \(f\in\mathrm{Fl}(\mathrm{H})\) . Indeed, the set of arrows \(f\) of G such that \(\varphi(f)=\varphi'(f)\) is the set of arrows of a subgroupoid of G, with vertex set S.

Let G be a groupoid with vertex set S, let X be a set, let \(p \colon \mathrm{X} \to \mathrm{S}\) be a map and let \(\varphi\) be an operation of G on X relative to \(p\) . Let \(\mathrm{G}_{\varphi}\) be the subquiver of \(p^{*}(\mathrm{G})\) whose set of vertices is X and whose set of arrows is the set of triples \((x, f, y) \in \mathrm{X} \times \mathrm{Fl}(\mathrm{G}) \times \mathrm{X}\) such that \(\varphi(f)(x) = y\) . It is a subgroupoid of \(p^{*}(\mathrm{G})\) .

Conversely, let \(\Gamma\) be a subgroupoid of \(p^{*}(\mathrm{G})\) ; suppose that the map \((x,f,y)\mapsto (x,f)\) from \(\mathrm{Fl}(\Gamma)\) to \(\mathrm{X}\times \mathrm{Fl}(\mathrm{G})\) is injective and that its image is the set of pairs \((x,f)\) such that \(p(x) = o(f)\) . There then exists a unique operation \(\varphi\) of the groupoid G on X such that one has \(\Gamma = \mathrm{G}_{\varphi}\) .

The orbits of the groupoid \(G_{\varphi}\) are called the orbits of the operation of G on X. By definition, they are the equivalence classes of the relation on X given by \(R\{x,y\}\) if and only if there exists \(f \in \mathrm{Fl}(G)\) such that \(\varphi(f)(x) = y\) .

Examples. --- 1) Let G be a group, let \(\mathcal{G}\) be the groupoid associated with G (II, p. 163, example 1). If X is a set, a (right) action of the groupoid \(\mathcal{G}\) on X is none other than a (right) action of the group G on X. Moreover, the orbits of the action of \(\mathcal{G}\) coincide with those of the group G.

\begin{enumerate}
\def\labelenumi{\arabic{enumi})}
\setcounter{enumi}{1}
\tightlist
\item
  Let \(\mathrm{G} = (\mathrm{S},\mathrm{F},o,t)\) be a groupoid. There exists a unique operation of \(\mathrm{G}\) on the set \(\mathrm{S}\) , relative to \(\mathrm{Id}_{\mathrm{S}}\) . It is given by \(\varphi(f)(a) = b\) if \(f \in \mathrm{G}_{a,b}\) . The orbits for this operation are the orbits of \(\mathrm{G}\) in the sense defined on p. 162.
\end{enumerate}

Let G be a groupoid, let S be the set of its vertices, let X be a set and let \(p: X \to S\) be a map. Let \(\varphi\) be an operation of G on X relative to p. One says that the groupoid G operates without monodromy on X if, for every point \(a \in S\) and every element \(f \in G_{a}\) , \(\varphi(f)\) is the identity of \(X_{a}\) . If the groupoid G is transitive, it suffices that this be so for one point of S.

Let \(a\) be a point of S and suppose that one has \(\varphi(f) = \mathrm{Id}_{\mathrm{X}_a}\) for all \(f \in \mathrm{G}_a\) . Let \(b\) be a point of S belonging to the orbit of \(a\) and let \(g\) be an arrow of G connecting \(a\) to \(b\) . For all \(f \in \mathrm{G}_b\) , \(\mathrm{Int}(g)(f) = gfg^{-1}\) is an element of \(\mathrm{G}_a\) ; one therefore has \(\mathrm{Id}_{\mathrm{X}_a} = \varphi(gfg^{-1}) = \varphi(g)\varphi(f)\varphi(g)^{-1}\) , so that \(\varphi(f) = \mathrm{Id}_{\mathrm{X}_b}\) . Let \(b\) be a point of S belonging to the orbit of \(a\) and let \(g\) , \(g'\) be arrows of G connecting \(a\) to \(b\) ; then \(g'g^{-1}\) is a loop at \(a\) , whence \(\varphi(g'g^{-1}) = \mathrm{Id}_{\mathrm{X}_a}\) and \(\varphi(g) = \varphi(g')\) .

